\section{Estrategia de estimación y muestra}

\subsection{Muestra de estimación}

La estimación se realiza sobre una muestra conjunta de viajes observados en los años 2024 y 2025. La unidad de observación es un viaje individual, para el cual se observa la alternativa de pago utilizada y se construyen atributos asociados a las alternativas disponibles. El conjunto de elección considerado en el modelo está compuesto por tres alternativas: \texttt{BIP}, \texttt{QR\_RED} y \texttt{QR\_OTHER}.

La muestra utilizada en la estimación principal corresponde a una submuestra estratificada del 2\% de la base interanual enriquecida. El submuestreo se realiza preservando la composición por partición temporal y por alternativa de pago observada. Es decir, dentro de cada combinación de período de datos y alternativa elegida se selecciona aleatoriamente la misma fracción de viajes, usando una semilla fija para asegurar reproducibilidad. De esta forma, la muestra conserva la presencia relativa de los distintos períodos y de las tres alternativas de pago: \texttt{BIP}, \texttt{QR\_RED} y \texttt{QR\_OTHER}.

Esta estrategia es relevante porque las alternativas de pago no tienen necesariamente la misma frecuencia en los datos. Un muestreo aleatorio simple podría alterar la proporción relativa de \texttt{BIP}, \texttt{QR\_RED} y \texttt{QR\_OTHER}, afectando la estabilidad de la estimación. Además, para evaluar que los resultados no dependan excesivamente del tamaño de muestra, se estimaron sensibilidades con muestras mayores en Larch. A diferencia de Biogeme en esta especificación, Larch permitió correr el mismo tipo de modelo con una mayor cantidad de datos sin generar problemas de memoria, por lo que se estimaron versiones al 5\% y al 10\% de la base. Estas corridas mostraron métricas de ajuste y patrones generales consistentes con la muestra de 2\%, lo que respalda el uso de la muestra reducida como base principal para la estimación y reporte en Biogeme.

En la corrida final reportada, la muestra contiene \(466{,}639\) observaciones y no se excluyen observaciones adicionales durante la estimación. El uso de una muestra interanual permite comparar patrones de elección entre 2024 y 2025, controlando explícitamente por el año de observación mediante la variable \texttt{ANIO\_2025}. De esta forma, el modelo no se estima separadamente por año, sino sobre una base pooled que permite capturar diferencias temporales sistemáticas dentro de una misma especificación.

La alternativa \texttt{BIP} se utiliza como referencia para las constantes específicas de alternativa. Por lo tanto, las constantes estimadas para \texttt{QR\_RED} y \texttt{QR\_OTHER} representan diferencias promedio de utilidad respecto de \texttt{BIP}, luego de controlar por el resto de variables incluidas en el modelo.

\subsection{Construcción de la base enriquecida}

La base de estimación se construye a partir de los viajes observados y se enriquece con información operacional, temporal y territorial. Primero, para cada viaje se incorporan variables alternativas-específicas de tiempo en vehículo, tiempo de espera inicial, tiempo de espera en transbordos y número de transbordos. Estas variables se calculan para cada alternativa de pago usando viajes comparables por origen, destino y tipo de pago, aplicando una corrección tipo \textit{leave-one-out} para evitar que el viaje observado contribuya mecánicamente a sus propios atributos.

En segundo lugar, se agregan variables de contexto temporal, demanda y oferta operacional. Las variables temporales identifican el año de observación y la franja horaria del viaje. La demanda local se calcula a partir del número de viajes observados en la zona de origen y franja temporal correspondiente. Las variables de oferta de transporte resumen la disponibilidad de paraderos de bus, servicios de bus y líneas de Metro en la zona de origen, combinando información espacial de paraderos, perfiles de oferta de bus y datos GTFS de Metro/Red.

En tercer lugar, cada viaje se asocia a una zona de origen \texttt{ZONA777}. Esta zona funciona como unidad espacial común para unir atributos territoriales provenientes de fuentes externas. El enriquecimiento territorial se realiza usando la zona donde comienza el viaje, no la zona de destino. Por lo tanto, las variables censales y OSM incluidas en el modelo describen el entorno residencial, urbano y de acceso al sistema de transporte en el punto de inicio del viaje.

Finalmente, se incorporan variables sociodemográficas del Censo 2024 y variables de OpenStreetMap. Las variables censales provienen de dos rutas: agregación espacial directa desde la base manzana-entidad hacia \texttt{ZONA777}, y espacialización de microdatos comunales hacia unidades \texttt{MANZENT} antes de agregarlos a \texttt{ZONA777}. Las variables OSM se construyen como densidades de objetos por kilómetro cuadrado dentro de cada zona. Todas las variables territoriales seleccionadas se estandarizan antes de ser incluidas en el modelo, lo que permite interpretar sus coeficientes como cambios asociados a un aumento de una desviación estándar en la característica territorial.

\subsection{Especificaciones estimadas}

La especificación central del reporte corresponde al modelo principal territorial. Este modelo combina atributos de viaje, controles temporales, variables de demanda y oferta operacional, variables sociodemográficas del Censo 2024 y variables de OpenStreetMap asociadas al entorno construido y a infraestructura de acceso al transporte. La motivación de esta especificación es evaluar si, además de los atributos operacionales del viaje, existen diferencias territoriales sistemáticas asociadas a la elección del medio de pago.

El modelo principal utiliza \texttt{BIP} como alternativa de referencia para las variables comunes al viaje o a la zona. Por lo tanto, para variables como año, franja horaria, demanda, oferta, Censo y OSM, los coeficientes de \texttt{QR\_RED} y \texttt{QR\_OTHER} se interpretan respecto de \texttt{BIP}. Esta normalización es directa desde el punto de vista de identificación del modelo y coincide con la forma usual de reportar un MNL con una alternativa base.

Además del modelo principal territorial, se estimaron sensibilidades para evaluar la robustez de la especificación. En particular, se compararon variantes con distintas combinaciones de variables OSM de infraestructura de acceso, incluyendo una especificación con accesos a Metro y otra con refugios de transporte más accesos a Metro. También se estimaron corridas en Larch con muestras de mayor tamaño para verificar que los patrones generales no dependieran exclusivamente de la muestra de 2\% usada en Biogeme.

\subsection{Criterios de evaluación}

Las especificaciones se evalúan combinando criterios estadísticos, computacionales e interpretativos. En primer lugar, se revisa la convergencia numérica del modelo. Una estimación se considera utilizable si el proceso de optimización converge sin errores, si el gradiente final es razonablemente bajo y si no aparecen problemas evidentes de identificación o parámetros inestables. Este criterio es especialmente relevante porque la especificación incluye un número elevado de parámetros y variables territoriales comunes a las alternativas.

En segundo lugar, se comparan medidas de ajuste global. La principal medida de ajuste es la log-verosimilitud final, donde valores menos negativos indican un mejor ajuste a los datos. También se reportan el criterio de información de Akaike (AIC) y el criterio de información bayesiano (BIC). Ambos penalizan la complejidad del modelo, pero BIC aplica una penalización más fuerte al número de parámetros, especialmente en muestras grandes. Por esta razón, una mejora en log-verosimilitud y AIC acompañada de un aumento leve en BIC no se interpreta automáticamente como evidencia en contra del modelo, sino como una señal de que el aporte incremental debe justificarse también por interpretabilidad y valor sustantivo.

En tercer lugar, se revisa la significancia estadística y estabilidad de los parámetros. Para los resultados principales se reportan errores estándar robustos, estadísticos \(t\) y valores \(p\). La interpretación se concentra no solo en si un parámetro es significativo, sino también en si su signo es estable, si su magnitud es razonable y si la lectura sustantiva es coherente con el fenómeno estudiado.

Finalmente, la selección del modelo principal no se basa únicamente en maximizar ajuste estadístico. Dado que el objetivo del reporte es describir un modelo interpretable de elección de medio de pago, se privilegia una especificación parsimoniosa, con variables territorialmente defendibles y con lectura útil para política pública. En particular, las variables de infraestructura de acceso se mantienen cuando aportan información sustantiva adicional sobre la experiencia de acceso al sistema, aun cuando su aporte en BIC sea más exigente de justificar.
